\unit{Exposé 149 — Duality Theorems for Coherent Algebraic Sheaves}
\sid{149}
\src{169}
\seminaren{(May 1957)}
\paperhead{DUALITY THEOREMS FOR COHERENT ALGEBRAIC SHEAVES}%
  {by Alexandre GROTHENDIECK}

The results that follow, inspired by Serre's \fquote{algebraic duality theorem},
were found in the winter of 1955 and the winter of 1956. They are established
very simply with the aid of fairly elementary results on the cohomology of
projective spaces [3], and through intensive use of the homological algebra of
Cartan-Eilenberg, in the form given in [2].

\fsection{\normalfont1}{\normalfont\underline{The} Ext's \underline{of sheaves of modules} ({[2]}, Chapters 3 and 4).}

Let $\symup{X}$ be a topological space equipped with a sheaf $\symup{\Osh}$ of rings with
identity (not necessarily commutative). One considers the abelian category
$\symup{\Ccat^{\Osh}}$ of sheaves of $\symup{\Osh}$-modules, also called $\symup{\Osh}$-Modules. Every
object of this category is known to admit an injective resolution, which makes
it possible to define in it the $\symup{\operatorname{Ext}}$ functors having the familiar
formal properties. More precisely, to avoid confusion, we denote by
$\symup{\operatorname{Hom}_{\Osh}(X;A,B)}$, or simply $\symup{\operatorname{Hom}(X;A,B)}$, the
abelian \underline{group} of $\symup{\Osh}$-homomorphisms from $\symup{A}$ to $\symup{B}$, whereas
$\symup{\Homsh_{\Osh}(A,B)}$ will denote the \underline{sheaf} of germs of homomorphisms from $\symup{A}$ to
$\symup{B}$ ($\symup{A,B\in\Ccat^{\Osh}}$). For fixed $\symup{A\in\Ccat^{\Osh}}$, one defines functors
$\symup{h_A}$ and $\symup{\hsh_A}$, taking values respectively in the category $\symup{\Ccat}$ of
abelian groups and the category $\symup{\Ccat^{\mathrm Z}}$ of sheaves of abelian groups on
$\symup{X}$, by the formulas
\begin{equation}\tag{1.1}
  \symup{h_A(B)=\operatorname{Hom}_{\Osh}(X;A,B)}
  \qquad
  \symup{\hsh_A(B)=\Homsh_{\Osh}(A,B)}
\end{equation}

The functors $\symup{h_A}$ and $\symup{\hsh_A}$ are covariant left-exact functors; one considers
their right derived functors, denoted respectively by
$\symup{\operatorname{Ext}_{\Osh}^{p}(X;A,B)}$ and $\symup{\Extsh_{\Osh}^{p}(A,B)}$. Thus, by
definition,
\begin{equation}\tag{1.2}
\left\{
\begin{aligned}
\symup{\operatorname{Ext}_{\Osh}^{p}(X;A,B)}
  &= \symup{(R^{p}h_A)(B)
   = H^{p}\bigl(\operatorname{Hom}_{\Osh}(X;A,C(B))\bigr)}\\
\symup{\Extsh_{\Osh}^{p}(A,B)}
  &= \symup{(R^{p}\hsh_A)(B)
   = H^{p}\bigl(\Homsh_{\Osh}(A,C(B))\bigr)}
\end{aligned}
\right.
\end{equation}

where the symbol $\symup{R^{p}}$ denotes passage to right derived functors, and where
$\symup{C(B)}$ denotes an arbitrary injective resolution of $\symup{B}$ in $\symup{\Ccat^{\Osh}}$. Let
$\symup{\Gamma:\Ccat^{\mathrm Z}\longrightarrow\Ccat}$
denote the \fquote{sections} functor. Recall that its right derived functors are
denoted by
$\symup{B\longrightarrow H^{p}(X,B)}$:

\src{170}
\begin{equation}\tag{1.3}
  \symup{H^{p}(X,B)=(R^{p}\Gamma)(B)=H^{p}\bigl(\Gamma(C(B))\bigr)}
\end{equation}
Clearly $\symup{h_A=\Gamma\hsh_A}$; on the other hand, one verifies that $\symup{\hsh_A}$
takes injective objects to $\symup{\Gamma}$-acyclic objects. The familiar conclusion is
the following.

\result{PROPOSITION}{1} \uline{For every} $\symup{\Osh}$\uline{-Module} $\symup{A}$, \uline{there exists a
cohomological spectral functor on} $\symup{\Ccat^{\Osh}}$ \uline{converging to the graded functor}
$\symup{\bigl(\operatorname{Ext}_{\Osh}^{*}(X;A,B)\bigr)}$, \uline{whose initial term is}
\begin{equation}\tag{1.4}
  \symup{E_{2}^{p,q}(A,B)=H^{p}\bigl(X,\Extsh_{\Osh}^{q}(A,B)\bigr)}
\end{equation}
From it one obtains \fquote{\uline{edge-homomorphisms}} and a five-term exact sequence, which we
shall not write down.

\result{COROLLARY}{1} \uline{If} $\symup{A}$ \uline{is locally isomorphic to} $\symup{\Osh^{n}}$, \uline{then there are
canonical isomorphisms}
\begin{equation}\tag{1.5}
  \symup{\operatorname{Ext}_{\Osh}^{p}(X;A,B)
  \overset{\sim}{\longleftarrow}
  H^{p}\bigl(X,\Homsh_{\Osh}(A,B)\bigr)}
\end{equation}
(given by edge-homomorphisms of the spectral sequence). \uline{In particular, there is
a canonical isomorphism}
\begin{equation}\tag{1.6}
  \symup{\operatorname{Ext}_{\Osh}^{p}(X;\Osh,B)=H^{p}(X,B)}
\end{equation}

To use these results, one must be able to describe the
$\symup{\Extsh_{\Osh}^{p}(A,B)}$. These functors are computed locally: if $\symup{U}$ is an open
subset of $\symup{X}$, then
\[
  \symup{\Extsh_{\Osh}^{p}(A,B)|U
  =\Extsh_{\Osh|U}^{p}(A|U,B|U)}
\]
as follows from the fact that the restriction to $\symup{U}$ of an injective
$\symup{\Osh}$-Module is an injective $\symup{(\Osh|U)}$-Module. Furthermore, for fixed
$\symup{x\in X}$, there are functorial homomorphisms
\begin{equation}\tag{1.7}
  \symup{\Homsh_{\Osh}(A,B)_x
  \longrightarrow
  \operatorname{Hom}_{\Osh_x}(A_x,B_x)}
\end{equation}
which extend uniquely to a homomorphism of cohomological functors in $\symup{B}$:
\begin{equation}\tag{1.8}
  \symup{\Extsh_{\Osh}^{p}(A,B)_x
  \longrightarrow
  \operatorname{Ext}_{\Osh_x}^{p}(A_x,B_x)}
\end{equation}

\result{PROPOSITION}{2} \uline{If} $\symup{A}$ \uline{is, in a neighborhood of} $\symup{x}$, \uline{isomorphic to the
cokernel of a homomorphism} $\symup{\Osh^{m}\longrightarrow\Osh^{n}}$, \uline{then} (1.7) \uline{is an isomorphism for every} $\symup{p}$. \uline{This holds in particular if}
$\symup{A}$ \uline{is a coherent} $\symup{\Osh}$-\uline{Module} [3].

\src{171}
\result{PROPOSITION}{3} \uline{Let} $\symup{L_{*}=(L_i)}$ \uline{be a left resolution of the}
$\symup{\Osh}$\uline{-Module} $\symup{A}$ \uline{by} $\symup{\Osh}$\uline{-Modules, each locally isomorphic to some} $\symup{\Osh^{n}}$.
\uline{Then} $\symup{\Extsh_{\Osh}(A,B)}$ \uline{identifies with}
$\symup{H^{*}\bigl(\Homsh_{\Osh}(L_{*},B)\bigr)}$, \uline{and}
$\symup{\operatorname{Ext}_{\Osh}(X;A,B)}$ \uline{identifies with the hypercohomology of}
$\symup{X}$ \uline{with respect to the complex} $\symup{\Homsh_{\Osh}(L_{*},B)}$.

The proof is standard: consider the bicomplex $\symup{\Homsh_{\Osh}(L_{*},C(B))}$,
where $\symup{C(B)}$ is an injective resolution of $\symup{B}$, and the natural homomorphisms
from $\symup{\Homsh_{\Osh}(L_{*},B)}$ and $\symup{\Homsh_{\Osh}(A,C(B))}$ into this bicomplex.

Finally, let us point out that the two $\symup{\operatorname{Ext}}$ functors introduced
in (1.2) are not only cohomological functors in $\symup{B}$, but \uline{cohomological bifunctors}, covariant in $\symup{B}$ and contravariant in $\symup{A}$.

\fsection{\normalfont2}{\normalfont\uline{Composition law in} Ext.}

The results of this section are due independently to CARTIER and YONEDA; see an
exposé by CARTIER [1] for details. Let $\symup{\underline{C}}$ be an abelian category, and let
$\symup{K}$ and $\symup{L}$ be two graded objects in $\symup{\underline{C}}$. One denotes by
$\symup{\operatorname{Hom}(K,L)}$ the graded abelian group whose component of degree
$\symup{n}$ consists of the homogeneous homomorphisms of degree $\symup{n}$ from $\symup{K}$ to $\symup{L}$
(that is, systems $\symup{(u_i)}$ of homomorphisms
$\symup{K^{i}\longrightarrow L^{i+n}}$).

If $\symup{K}$ and $\symup{L}$ are complexes (with differential operators of degree $\symup{+1}$, to
fix ideas), one equips $\symup{\operatorname{Hom}(K,L)}$ with the differential operator
given by
\begin{equation}\tag{2.1}
  \symup{\delta(u)=du+(-1)^{n+1}ud,\qquad\text{where }n=\deg(u)}
\end{equation}
making it a complex with differential of degree $\symup{+1}$. The cycles of degree $\symup{n}$
are the operators of degree $\symup{n}$ that commute with $\symup{d}$ (as homogeneous
operators). One may therefore consider
$\symup{H^{*}\bigl(\operatorname{Hom}(K,L)\bigr)}$; it is an invariant of the homotopy
types of $\symup{K}$ and $\symup{L}$, and may be denoted by $\symup{H^{*}(K,L)}$.

Given a third complex $\symup{M}$, composition of homomorphisms defines a pairing
$\symup{\operatorname{Hom}(K,L)\times\operatorname{Hom}(L,M)
\longrightarrow\operatorname{Hom}(K,M)}$
compatible with the differentials, hence, after passing to cohomology, pairings
\begin{equation}\tag{2.2}
  \symup{H^{*}(K,L)\times H^{*}(L,M)\longrightarrow H^{*}(K,M)}
\end{equation}
denoted $\symup{(u,v)\longrightarrow vu}$. These pairings satisfy the evident associativity
property. In particular, $\symup{H^{*}(K,K)}$ is an associative graded ring with
identity, and $\symup{H^{*}(K,L)}$ (resp. $\symup{H^{*}(L,K)}$) is a graded right
(resp. left) module over this ring, and so on. In degree $\symup{0}$, (2.2)
reduces to composition of permitted homomorphisms of complexes. Finally, an exact sequence of
complexes
\src{172}
$\symup{0\longrightarrow K'\longrightarrow K\longrightarrow K''\longrightarrow 0}$, such that $\symup{K'^{i}}$ is a direct summand of $\symup{K^{i}}$ for every $\symup{i}$, gives rise to an
exact sequence of complexes of groups $\symup{\operatorname{Hom}(K'',L)}$, and so on,
and hence to a coboundary operator
$\symup{H^{i}(K',L)\longrightarrow H^{i+1}(K'',L)}$.
The coboundary operators associated with an exact sequence in $\symup{L}$ are defined
similarly. The pairings (2.2) are compatible, in the usual sense, with these
coboundary operators.

Now suppose that $\symup{C}$ is a category such that every object $\symup{A}$ of $\symup{C}$ admits an
injective resolution $\symup{C(A)}$. One then finds, using one of the many variants of
the bicomplex theorem, that
\[
  \symup{H^{*}\bigl(C(A),C(B)\bigr)
  =H^{*}\bigl(\operatorname{Hom}(C(A),C(B))\bigr)}
\]
is canonically isomorphic to
\[
  \symup{H^{*}\bigl(\operatorname{Hom}(A,C(B))\bigr)
  =\operatorname{Ext}^{*}(A,B).}
\]
The coboundary operators considered above give the coboundary homomorphisms for
$\symup{\operatorname{Ext}}$. Moreover, the pairings (2.2) here give associative
pairings compatible with the coboundary homomorphisms:
\begin{equation}\tag{2.3}
  \symup{\operatorname{Ext}^{*}(A,B)\times\operatorname{Ext}^{*}(B,C)
  \longrightarrow\operatorname{Ext}^{*}(A,C)}
\end{equation}
In particular, $\symup{\operatorname{Ext}^{*}(A,A)}$ is an associative graded ring with
identity, and so on. (One proves similarly that the $\symup{\operatorname{Ext}}$
functors operate on the derived functors of an arbitrary functor; we shall not
use this fact here.)

When the category in question is the category $\symup{\Ccat^{\Osh}}$ of
$\symup{\Osh}$-Modules on $\symup{X}$, one therefore obtains pairings
\begin{equation}\tag{2.4}
  \symup{\operatorname{Ext}_{\Osh}^{p}(X;A,B)
  \times\operatorname{Ext}_{\Osh}^{q}(X;B,C)
  \longrightarrow\operatorname{Ext}_{\Osh}^{p+q}(X;A,C)}
\end{equation}
which may be made explicit as indicated above. Moreover, the same developments,
with the category of abelian groups replaced by the category of sheaves of
abelian groups on $\symup{X}$ and the functors $\symup{\operatorname{Hom}}$ replaced by the
functors $\symup{\Homsh}$, also define pairings having the same formal properties and,
this time, of a \fquote{local nature}:
\begin{equation}\tag{2.5}
  \symup{\Extsh_{\Osh}^{p}(A,B)\times\Extsh_{\Osh}^{q}(B,C)
  \longrightarrow\Extsh_{\Osh}^{p+q}(A,C)}
\end{equation}
These are sharpened by observing that the homomorphisms (1.8) are compatible
with the pairings between the $\symup{\operatorname{Ext}}$.

\src{173}
Finally, recall that there is also a multiplicative structure among the functors
$\symup{H^{p}(X,A)}$ (the cup product). One then finds that the spectral sequences of
Proposition 1 are compatible with the multiplicative structures; more
precisely, one has a pairing of the spectral sequences $\symup{E(A,B)}$ and $\symup{E(B,C)}$
into the spectral sequence $\symup{E(A,C)}$, which at the abutment reduces to the pairing
between global $\symup{\operatorname{Ext}}$, and on the initial term to the
pairing deduced, in the right-hand side of (1.4), from the cup product and the
pairings of local $\symup{\operatorname{Ext}}$. In particular, it follows that
the \fquote{edge-homomorphisms}
\begin{equation}\tag{2.6}
  \symup{\operatorname{Ext}_{\Osh}^{n}(X;A,B)
  \longrightarrow H^{0}\bigl(X;\Extsh_{\Osh}^{n}(A,B)\bigr)}
\end{equation}
\begin{equation}\tag{2.7}
  \symup{H^{n}\bigl(X,\Homsh_{\Osh}(A,B)\bigr)
  \longrightarrow\operatorname{Ext}_{\Osh}^{n}(X;A,B)}
\end{equation}
are compatible with the multiplicative structures. Thus, if one restricts to
sheaves locally isomorphic to some $\symup{\Osh^{m}}$, this completely describes
composition between global $\symup{\operatorname{Ext}}$ in terms of the cup
product, in view of the isomorphism (1.5).

\fsection{\normalfont3}{\normalfont\uline{Results in local cohomology}.}

Let $\symup{A}$ be a commutative ring with identity, equipped with an ideal $\symup{J}$. We shall define, for a
varying $\symup{A}$-module $\symup{M}$, functorial homomorphisms
\begin{equation}\tag{3.1}
\left\{
\begin{aligned}\symup{\operatorname{Ext}_{A}^{p}(A/J,M)}&\symup{\longrightarrow
  \operatorname{Hom}_{A}\!\left(\bigwedge^{p}J/J^{2},M\otimes A/J\right)}\\\symup{\operatorname{Tor}_{p}^{A}(A/J,M)}&\symup{\longleftarrow
  \left(\bigwedge^{p}J/J^{2}\right)
  \otimes\operatorname{Hom}_{A}(A/J,M)}\end{aligned}
\right.
\end{equation}
where the tensor and exterior products are taken over the ring $\symup{A}$; notice also
that $\symup{J/J^{2}}$ is in fact an $\symup{A/J}$-module, and that its exterior powers as an
$\symup{A}$-module or as an $\symup{A/J}$-module are the same. Defining the homomorphisms (3.1)
amounts to defining, for every system
$\symup{\underline{x}=(x_1,\ldots,x_p)}$ of elements of $\symup{J}$, homomorphisms
$\symup{\varphi_{\underline{x}}}$:
\begin{equation}\tag{3.2}
\left\{
\begin{aligned}\symup{\varphi_{\underline{x}}:
  \operatorname{Ext}_{A}^{p}(A/J,M)}&\symup{\longrightarrow M\otimes A/J}\\\symup{\varphi_{\underline{x}}:
  \operatorname{Hom}_{A}(A/J,M)}&\symup{\longrightarrow
  \operatorname{Tor}_{p}^{A}(A/J,M)}\end{aligned}
\right.
\end{equation}
in such a way that the following conditions are satisfied:

\src{174}
\begin{enumerate}
\item[\textup{i.}] $\symup{\varphi_{x_1,\ldots,x_p}}$ depends in an alternating
$\symup{A}$-multilinear fashion on the system of $\symup{x_i\in J}$
\item[\textup{ii.}] $\symup{\varphi_{x_1,\ldots,x_p}}$ is zero whenever one of the
$\symup{x_i}$ lies in $\symup{J^{2}}$.
\end{enumerate}
In fact, ii. follows from i., since $\symup{a\varphi_{\underline{x}}=0}$ for $\symup{a\in J}$, as
is seen by observing that all the modules in (3.2) are annihilated by $\symup{J}$.

To define $\symup{\varphi_{\underline{x}}}$, consider the complex
$\symup{K_{\underline{x}}}$ whose underlying $\symup{A}$-module is $\symup{\bigwedge A^{p}}$ and whose
differential is contraction $\symup{i_{\underline{x}}}$ by $\symup{\underline{x}}$, regarded
as the linear form on $\symup{A^{p}}$ with components $\symup{x_1,\ldots,x_p}$. The differential
has degree $\symup{-1}$, the degrees are nonnegative, and the cohomology of this complex
in dimension $\symup{0}$ is $\symup{A/(x_1A+\cdots+x_pA)}$. Since the $\symup{x_i}$ belong to $\symup{J}$, this
gives an augmentation $\symup{K_{\underline{x},0}\longrightarrow A/J}$. Thus
$\symup{K_{\underline{x}}}$ appears as an augmented \uline{free} complex with augmentation
module $\symup{A/J}$. In the familiar way, one obtains homomorphisms
\[
\symup{\operatorname{Ext}_{A}^{*}\bigl(H_0(K_{\underline{x}}),M\bigr)
\longrightarrow
H^{*}\bigl(\operatorname{Hom}_{A}(K_{\underline{x}},M)\bigr)\quad\text{and}\quad
\operatorname{Tor}_{*}^{A}\bigl(H_0(K_{\underline{x}}),M\bigr)
\longleftarrow
H_{*}\bigl(K_{\underline{x}}\otimes M\bigr)}
\]
whence, by composing with the homomorphisms on $\symup{\operatorname{Ext}}$ and
$\symup{\operatorname{Tor}}$ induced by the augmentation homomorphism
$\symup{H_0(K_{\underline{x}})\longrightarrow A/J}$, one obtains homomorphisms
\begin{equation}\tag{3.3}
\left\{
\begin{aligned}\symup{\psi_{\underline{x}}:
\operatorname{Ext}_{A}^{*}(A/J,M)}&\symup{\longrightarrow H^{*}\bigl(\operatorname{Hom}_{A}(K_{\underline{x}},M)\bigr)}\\\symup{\psi_{\underline{x}}:
\operatorname{Tor}_{*}^{A}(A/J,M)}&\symup{\longleftarrow H_{*}\bigl(K_{\underline{x}}\otimes M\bigr)}\end{aligned}
\right.
\end{equation}
One sees immediately that in the top dimension $\symup{p}$, the cohomology of the
right-hand sides is $\symup{M/(x_1M+\cdots+x_pM)}$ (resp. the set of elements of
$\symup{M}$ annihilated by every $\symup{x_i}$). Since the $\symup{x_i}$ lie in $\symup{J}$, one obtains
homomorphisms
\begin{equation}\tag{3.4}
\left\{
\begin{aligned}\symup{H^{p}\bigl(\operatorname{Hom}_{A}(K_{\underline{x}},M)\bigr)}&\symup{\longrightarrow M\otimes A/J}\\\symup{H_{p}\bigl(K_{\underline{x}}\otimes M\bigr)}&\symup{\longleftarrow\operatorname{Hom}_{A}(A/J,M)}\end{aligned}
\right.
\end{equation}
Composing the homomorphisms (3.3) and (3.4) gives the homomorphisms (3.2) that we
wished to define. Verifying i. is tedious but presents no difficulty.

\src{175}
\result{PROPOSITION}{4} \uline{Let} $\symup{A}$ \uline{be a commutative ring with identity, and let}
$\symup{(x_1,\ldots,x_p)}$ \uline{be a sequence of elements of} $\symup{A}$ \uline{such that, for}
$\symup{1\leq i\leq p}$, \uline{the image of} $\symup{x_i}$ \uline{in the quotient of} $\symup{A}$ \uline{by the ideal
generated by} $\symup{(x_1,\ldots,x_{i-1})}$ \uline{is not a divisor of} $\symup{0}$. \uline{Let} $\symup{J}$ \uline{be the ideal
generated by the} $\symup{x_i}$. \uline{Then} $\symup{J/J^{2}}$ \uline{is a free} $\symup{(A/J)}$\uline{-module with basis the
canonical images of the} $\symup{x_i}$, \uline{the complex} $\symup{K_{\underline{x}}}$ \uline{is a free
resolution of} $\symup{A/J}$, \uline{and for every} $\symup{A}$\uline{-module} $\symup{M}$, \uline{the homomorphisms} (3.1) \uline{in
degree} $\symup{p}$ \uline{are bijective. The same is true of the analogous homomorphisms
defined in arbitrary degrees} $\symup{i}$, \uline{provided that} $\symup{J.M=0}$.

(The essential point, from which all the others follow, is the acyclicity of
$\symup{K_{\underline{x}}}$, a well-known fact under the stated hypotheses.)

\result{COROLLARY}{1} \uline{With} $\symup{A}$ \uline{and} $\symup{J}$ \uline{as above, suppose in addition that} $\symup{A}$
\uline{is a regular affine algebra of} $\symup{\dim n}$ \uline{over a perfect field} $\symup{\underline{k}}$, \uline{and that}
$\symup{A/J}$ \uline{is a regular affine algebra. Denote by} $\symup{\Omega^{1}(A)}$ \uline{and}
$\symup{\Omega^{1}(A/J)}$ \uline{the modules of Kähler differentials. Then there is a canonical
isomorphism compatible with localization},
\begin{equation}\tag{3.5}
  \symup{\operatorname{Ext}_{A}^{p}\bigl(\Omega^{n-p}(A/J),\Omega^{n}(A)\bigr)=A/J}
\end{equation}
Indeed, $\symup{\Omega^{n-p}(A/J)}$ is a projective $\symup{(A/J)}$-module of rank $\symup{1}$, and likewise
$\symup{\Omega^{n}(A)}$ is a projective $\symup{A}$-module of rank $\symup{1}$, so the left-hand
side equals
\[
  \symup{\operatorname{Ext}_{A}^{p}(A/J,A)
  \otimes\Omega^{n-p}(A/J)'\otimes\Omega^{n}(A)}
\]
(where the symbol \fquote{$\symup{\prime}$} denotes the dual $\symup{(A/J)}$-module). The tensor product of the
last two factors identifies with $\symup{\bigwedge^{p}(J/J^{2})}$, so the whole expression
identifies with
$\symup{\operatorname{Ext}_{A}^{p}\!\left(A/J,\bigwedge^{p}(J/J^{2})\right)}$,
and therefore, by the proposition, with
\[
  \symup{\operatorname{Hom}_{A}\!\left(\bigwedge^{p}J/J^{2},
  \bigwedge^{p}J/J^{2}\right),}
\]
i.e. with $\symup{A/J}$. In particular,
$\symup{\operatorname{Ext}_{A}^{p}\bigl(\Omega^{n-p}(A/J),\Omega^{n}(A)\bigr)}$ contains
a distinguished element corresponding to the identity element of $\symup{A/J}$, called
the \uline{fundamental class} of the ideal $\symup{J}$ in $\symup{A}$. (It can in fact be defined
under considerably broader hypotheses.) Corollary 1 can be put in a more
geometric and more global form:

\result{COROLLARY}{2} \uline{Let} $\symup{X}$ \uline{be a nonsingular variety over an algebraically
closed field} $\symup{\underline{k}}$, \uline{let} $\symup{Y}$ \uline{be a closed nonsingular subvariety of} $\symup{X}$, \uline{let}
$\symup{\Osh_X}$ \uline{be the sheaf of local rings of} $\symup{X}$, \uline{and let} $\symup{\Osh_Y}$ \uline{be the sheaf of
local rings of} $\symup{Y}$, \uline{regarded as a quotient sheaf of} $\symup{\Osh_X}$. \uline{Let} $\symup{n}$ \uline{be the
dimension of} $\symup{X}$ \uline{and} $\symup{n-p}$ \uline{that of} $\symup{Y}$.

\src{176}
\uline{Let} $\symup{\underline{\Omega}_X}$ (resp. $\symup{\underline{\Omega}_Y}$) \uline{be the sheaf of germs of regular
differential forms on} $\symup{X}$ (resp. $\symup{Y}$). \uline{Then there are canonical
isomorphisms}:
\begin{equation}\tag{3.6}
  \symup{\Extsh_{\Osh_X}^{p}\bigl(\underline{\Omega}_Y^{n-p},\underline{\Omega}_X^{n}\bigr)=\Osh_Y}
\end{equation}
\uline{or equivalently}
\begin{equation}\tag{3.6 bis}
  \symup{\Extsh_{\Osh_X}^{p}\bigl(\Osh_Y,\underline{\Omega}_X^{n}\bigr)=\underline{\Omega}_Y^{n-p}}
\end{equation}
Formula (3.6 bis) can serve as the \uline{definition} of $\symup{\underline{\Omega}_Y^{n-p}}$ when $\symup{Y}$ is a
singular variety. More precisely:

\result{PROPOSITION}{5} \uline{Let} $\symup{Y}$ \uline{be an algebraic subset of dimension} $\symup{q=n-p}$ \uline{in
a nonsingular algebraic variety} $\symup{X}$ \uline{of dimension} $\symup{n}$. \uline{Let} $\symup{F}$ \uline{be a coherent
algebraic sheaf on} $\symup{X}$ \uline{whose support is contained in} $\symup{Y}$, \uline{and let} $\symup{L}$ \uline{be a
locally free algebraic sheaf on} $\symup{X}$. \uline{Then the sheaves}
$\symup{\Extsh_{\Osh_X}^{i}(F,L)}$ \uline{vanish for} $\symup{i<p}$, \uline{while for} $\symup{i=p}$ \uline{there is a canonical
isomorphism}
\begin{equation}\tag{3.7}
  \symup{\Extsh_{\Osh_X}^{p}(F,L)
  =\Homsh_{\Osh_X}\!\left(F,
  \Extsh^{p}(\Osh_X/J,L)\right)}
\end{equation}
\uline{where} $\symup{J}$ \uline{denotes an arbitrary sheaf of ideals on} $\symup{X}$ \uline{that annihilates} $\symup{F}$ \uline{and
whose zero set is} $\symup{Y}$. \uline{In particular, if} $\symup{F}$ \uline{is a coherent algebraic sheaf on}
$\symup{Y}$, \uline{then}
\begin{equation}\tag{3.7 bis}
  \symup{\Extsh_{\Osh_X}^{p}(F,L)
  =\Homsh_{\Osh_Y}\!\left(F,\Extsh^{p}(\Osh_Y,L)\right).}
\end{equation}
\uline{Finally, still taking} $\symup{F}$ \uline{to be a coherent algebraic sheaf on} $\symup{Y}$, \uline{the sheaves}
$\symup{E^{i}(F)=\Extsh_{\Osh_X}^{p+i}(F,\underline{\Omega}_X^{n})}$
\uline{do not depend on the way of immersing the algebraic space} $\symup{Y}$ \uline{into a
nonsingular variety} $\symup{X}$.

Since the question is local, one may suppose that $\symup{X}$ is affine and
$\symup{L=\Osh_X}$. This reduces matters to a question in commutative algebra, indeed in
local algebra: if $\symup{A}$ is a regular local ring and $\symup{M}$ an $\symup{A}$-module whose
support has dimension $\symup{\leq q=n-p}$, one must prove that
$\symup{\operatorname{Ext}_{A}^{i}(M,A)=0}$ for $\symup{i<p}$, and that
$\symup{\operatorname{Ext}_{A}^{p}(M,A)
=\operatorname{Hom}_{A}\!\left(M,\operatorname{Ext}^{p}(A/J,A)\right)}$,
where $\symup{J}$ is any ideal of \fquote{dimension} $\symup{\leq q}$ that annihilates $\symup{M}$. For the
first assertion, one proceeds by induction on $\symup{q}$: an immediate dévissage reduces to
the case where $\symup{M}$ is of the form $\symup{A/J}$; replacing $\symup{J}$ by a smaller ideal, using the induction
hypothesis and the exact sequence of $\symup{\operatorname{Ext}}$, then reduces
to the case where $\symup{J}$ is generated by a \fquote{system of parameters} as in
Proposition 4, where the result is immediate. The preceding result implies that
if $\symup{J}$ is an ideal

\src{177}
fixed of \fquote{dimension} $\symup{\leq q}$, then the contravariant functor
$\symup{E(M)=\operatorname{Ext}_{A}^{p}(M,A)}$ on the category of $\symup{A/J}$-modules is left
exact; moreover, it takes direct sums to direct products, from which it follows
readily that
$\symup{E(M)=\operatorname{Hom}_{A}\bigl(M,E(A/J)\bigr)}$.
Finally, the last assertion of Proposition 5 is subtler and follows from an
intrinsic characterization of the $\symup{E^{i}(F)}$ by means of a local duality theorem
that cannot be given here.

\result{COROLLARY}{} \uline{Denote by} $\symup{\underline{\omega}_Y^{q}}$ \uline{the sheaf}
$\symup{\Extsh_{\Osh_X}^{p}(\Osh_Y,\underline{\Omega}_X^{n})}$. \uline{Then there is a functorial isomorphism for coherent algebraic sheaves}
$\symup{F}$ \uline{on} $\symup{Y}$:
\begin{equation}\tag{3.8}
  \symup{\Extsh_{\Osh_X}^{p}(F,\underline{\Omega}_X^{n})
  =\Homsh_{\Osh_X}(F,\underline{\omega}_Y^{q})}
\end{equation}

\fsection{4}{\uline{Cohomology class associated with a subvariety.}}

Throughout what follows, $\symup{X}$ denotes an algebraic set of dimension $\symup{n}$, defined
over a field $\symup{\underline{k}}$, which we suppose algebraically closed for simplicity. Except
in Section 6, $\symup{X}$ is assumed nonsingular. One denotes by $\symup{\Osh_X}$ the sheaf of
local rings on $\symup{X}$, and by
$\symup{\underline{\Omega}_X^{*}=\bigoplus\limits_p\underline{\Omega}_X^{p}}$
the sheaf of germs of differential forms on $\symup{X}$. If $\symup{Y}$ is a closed subset of $\symup{X}$, coherent algebraic sheaves on
$\symup{Y}$ are identified with coherent sheaves on $\symup{X}$ that vanish off $\symup{Y}$; this
applies in particular to $\symup{\Osh_Y}$ and $\symup{\underline{\Omega}_Y}$.

\result{LEMMA}{1} \uline{Let} $\symup{F}$ \uline{be a coherent algebraic sheaf on} $\symup{X}$ \uline{whose support
has dimension} $\symup{\leq n-p}$, \uline{and let} $\symup{L}$ \uline{be a locally free coherent algebraic sheaf
on} $\symup{X}$. \uline{Then} $\symup{\operatorname{Ext}_{\Osh_X}^{i}(X;F,L)}$ \uline{vanishes for} $\symup{i<p}$, \uline{and
there is a canonical isomorphism}
\begin{equation}\tag{4.1}
  \symup{\operatorname{Ext}_{\Osh}^{p}(X;F,L)
  =H^{0}\bigl(X,\Extsh_{\Osh}^{p}(F,L)\bigr)}
\end{equation}
\uline{If} $\symup{F}$ \uline{is a coherent algebraic sheaf on a closed subset} $\symup{Y}$ \uline{of} $\symup{X}$ \uline{of dimension}
$\symup{\leq n-p}$, \uline{there is a canonical isomorphism}
\begin{equation}\tag{4.1 bis}
  \symup{\Extsh_{\Osh_X}^{p}(F,L)
  =\Homsh_{\Osh_X}\!\left(
    F\otimes L'\otimes\underline{\Omega}_X^{n},\underline{\omega}_Y^{n-p}
  \right)}
\end{equation}
\uline{where} $\symup{\underline{\omega}_Y^{n-p}}$ \uline{is the sheaf on} $\symup{Y}$ \uline{defined in the corollary to
Proposition 5 (and identifies with} $\symup{\underline{\Omega}_Y^{n-p}}$ \uline{if} $\symup{Y}$ \uline{is nonsingular}).

Formula (4.1) is an immediate consequence of the spectral sequence of
Proposition 1 and Proposition 5; by formula (3.8), one may write

\src{178}
\[
\symup{\Extsh_{\Osh_X}^{p}(F,L)
=L\otimes(\underline{\Omega}_X^{n})'\otimes
  \Extsh_{\Osh_X}^{p}(F,\underline{\Omega}_X^{n})
=L\otimes(\underline{\Omega}_X^{n})'\otimes
  \Homsh_{\Osh_X}(F,\underline{\omega}_Y^{q})
=\Homsh_{\Osh_X}\!\left(
  F\otimes L'\otimes\underline{\Omega}_X^{n},\underline{\omega}_Y^{q}\right)}
\]
where $\symup{q=n-p}$, which gives formula (4.1 bis).

Let us take, in particular, $\symup{F=\Osh_Y}$, $\symup{L=\underline{\Omega}_X^{p}}$; taking into account that
$\symup{\underline{\Omega}_X^{n}\otimes(\underline{\Omega}_X^{p})'=\underline{\Omega}_X^{n-p}}$, one obtains a canonical
isomorphism
\begin{equation}\tag{4.2}
  \symup{\operatorname{Ext}_{\Osh_X}^{p}(X;\Osh_Y,\underline{\Omega}_X^{p})
  =\operatorname{Hom}_{\Osh_X}
  (X;\underline{\Omega}_X^{n-p},\underline{\omega}_Y^{n-p})}
\end{equation}
Let us now suppose $\symup{Y}$ \uline{nonsingular} for simplicity, so that
$\symup{\underline{\omega}_Y^{n-p}=\underline{\Omega}_Y^{n-p}}$. There is a natural homomorphism from
$\symup{\underline{\Omega}_X^{n-p}}$ to $\symup{\underline{\Omega}_Y^{n-p}}$, hence a canonical section $\symup{s_Y}$ of the
sheaf $\symup{\Extsh_{\Osh_X}^{p}(\Osh_Y,\underline{\Omega}_X^{p})}$, which we shall call, if all components of $\symup{Y}$ have
dimension $\symup{n-p}$, the \uline{fundamental section} of the sheaf
$\symup{\Extsh_{\Osh_X}^{p}(\Osh_Y,\underline{\Omega}_X^{p})}$. By (4.1), this section defines an
element of $\symup{\operatorname{Ext}_{\Osh_X}^{p}(X;\Osh_Y,\underline{\Omega}_X^{p})}$. The natural
homomorphism $\symup{\Osh_X\longrightarrow\Osh_Y}$ in turn defines a homomorphism
\[
  \symup{\operatorname{Ext}_{\Osh_X}^{p}(X;\Osh_Y,\underline{\Omega}_X^{p})
  \longrightarrow
  \operatorname{Ext}_{\Osh_X}^{p}(X;\Osh_X,\underline{\Omega}_X^{p})
  =H^{p}(X,\underline{\Omega}_X^{p}).}
\]
Thus one obtains an element of $\symup{H^{p}(X,\underline{\Omega}_X^{p})}$, denoted $\symup{P_X(Y)}$ and
called the \uline{cohomology class of} $\symup{Y}$ \uline{in} $\symup{X}$. It is therefore obtained from
the section $\symup{s_Y}$ of $\symup{\Extsh_{\Osh_X}^{p}(\Osh_Y,\underline{\Omega}_X^{p})}$ through the
following diagram of homomorphisms:
\begin{equation}\tag{4.3}
\symup{H^{p}(X,\underline{\Omega}_X^{p})
=\operatorname{Ext}_{\Osh_X}^{p}(X;\Osh_X,\underline{\Omega}_X^{p})
\longleftarrow\operatorname{Ext}_{\Osh_X}^{p}(X;\Osh_Y,\underline{\Omega}_X^{p})
\xrightarrow{\ \sim\ }
H^{0}\!\left(X,\Extsh_{\Osh_X}^{p}(\Osh_Y,\underline{\Omega}_X^{p})\right)}
\end{equation}

Let us call a \uline{nonsingular cycle} of dimension $\symup{n-p}$ an element of the free
abelian group generated by the irreducible nonsingular subvarieties of dimension
$\symup{n-p}$ in $\symup{X}$. Then the function $\symup{Y\longrightarrow P(Y)}$ extends to a homomorphism
from the group of nonsingular cycles of dimension $\symup{n-p}$ on $\symup{X}$ to the group
$\symup{H^{p}(X,\underline{\Omega}_X^{p})}$.

Let $\symup{Z^{n-p}}$ and $\symup{Z'^{n-p'}}$ be nonsingular cycles of dimensions $\symup{n-p}$ and
$\symup{n-p'}$; they are said to \uline{intersect transversely}, if every component of
$\symup{Z}$ intersects every component of $\symup{Z'}$ transversely. The cycle $\symup{Z.Z'}$ is then
defined; it is a nonsingular cycle of dimension $\symup{n-p-p'}$. With this understood,
one has the

\src{179}
\result{THEOREM}{1} \uline{If} $\symup{Z^{n-p}}$ \uline{and} $\symup{Z'^{n-p'}}$ \uline{are nonsingular cycles that
intersect transversely, then one has}
\begin{equation}\tag{4.4}
  \symup{P_X(Z.Z')=P_X(Z).P_X(Z')}
\end{equation}
\uline{where the product on the right is the cup product}
\[
  \symup{H^{p}(X,\underline{\Omega}_X^{p})\times H^{p'}(X,\underline{\Omega}_X^{p'})
  \longrightarrow H^{p+p'}(X,\underline{\Omega}_X^{p+p'}).}
\]
(One assumes that $\symup{X}$ is isomorphic to a locally closed subset of a projective
space)

We use this last hypothesis only to be able to conclude that every coherent algebraic
sheaf on $\symup{X}$ is a quotient of a locally free coherent algebraic sheaf (SERRE),
and hence admits a left resolution by locally free sheaves. To prove Theorem 1,
one may suppose that $\symup{Z}$ and $\symup{Z'}$ are irreducible nonsingular subvarieties $\symup{Y}$
and $\symup{Y'}$ meeting transversely. Let $\symup{L_{*}}$ be a left resolution of $\symup{\Osh_Y}$ by
locally free sheaves; then, by Proposition 3, the diagram of homomorphisms (4.3)
identifies with the diagram
\[
\symup{R^{p}\Gamma(\underline{\Omega}_X^{p})
\xleftarrow{\alpha}
(\underline{\underline{R}}^{p}\Gamma)\!\left(\Homsh_{\Osh_X}(L_{*},\underline{\Omega}_X^{p})\right)
\xrightarrow{\overset{\beta}{\sim}}
\Gamma\!\left(H^{p}\bigl(\Homsh_{\Osh_X}(L_{*},\underline{\Omega}_X^{p})\bigr)\right)}
\]
where $\symup{\Gamma}$ is the \fquote{group of sections} functor on the category of sheaves of
abelian groups on $\symup{X}$, $\symup{\underline{\underline{R}}^{p}\Gamma}$ denotes its hypercohomology in
dimension $\symup{p}$, and $\symup{R^{p}\Gamma}$ its $\symup{p}$-th derived functor. For simplicity, one
will suppose that $\symup{L_0=\Osh_X}$ and that the augmentation
$\symup{L_0\longrightarrow\Osh_Y}$ is the natural homomorphism (which is permissible);
then $\symup{\alpha}$ is induced by the homomorphism of complexes
$\symup{\Osh_X\longrightarrow L}$ (with $\symup{\Osh_X}$ regarded as a complex concentrated
in degree $\symup{0}$), taking account of the equality
$\symup{\underline{\underline{R}}^{p}\Gamma(K)=R^{p}\Gamma(K_0)}$ when $\symup{K}$ is a complex of sheaves
concentrated in degree $\symup{0}$. The homomorphism $\symup{\beta}$ is a familiar \fquote{edge-homomorphism}. Let us consider an analogous diagram for a locally free resolution
$\symup{L'_{*}}$ of $\symup{\Osh_{Y'}}$, and let us consider the commutative diagram of pairings:
\begin{equation}\tag{4.5}
\resizebox{\textwidth}{!}{$
\left\{\begin{array}{ccc}\symup{R^{p}\Gamma(\underline{\Omega}_X^{p})}&\symup{\longleftarrow
\underline{\underline{R}}^{p}\!\Gamma\left(\Homsh_{\Osh_X}(L_{*},\underline{\Omega}_X^{p})\right)}&\symup{\xrightarrow{\sim}
\Gamma\bigl(H^{p}\!\left(\Homsh_{\Osh_X}(L_{*},\underline{\Omega}_X^{p})\right)\bigr)}\\[3pt]\symup{\times}&\symup{\times}&\symup{\times}\\[3pt]\symup{R^{p'}\Gamma(\underline{\Omega}_X^{p'})}&\symup{\longleftarrow
\underline{\underline{R}}^{p'}\Gamma\!\left(\Homsh_{\Osh_X}(L'_{*},\underline{\Omega}_X^{p'})\right)}&\symup{\xrightarrow{\sim}
\Gamma\bigl(H^{p'}\!\left(\Homsh_{\Osh_X}(L'_{*},\underline{\Omega}_X^{p'})\right)\bigr)}\\[3pt]\symup{\downarrow}&\symup{\downarrow}&\symup{\downarrow}\\[3pt]\symup{R^{p+p'}\Gamma(\underline{\Omega}_X^{p+p'})}&\symup{\longleftarrow
\underline{\underline{R}}^{p+p'}\Gamma\!\left(
\Homsh_{\Osh_X}(L_{*}\otimes L'_{*},\underline{\Omega}_X^{p+p'})\right)}&\symup{\xrightarrow{\sim}
\Gamma\bigl(H^{p+p'}\!\left(
\Homsh_{\Osh_X}(L_{*}\otimes L'_{*},\underline{\Omega}_X^{p+p'})\right)\bigr)}\end{array}\right.
$}
\end{equation}

\src{180}
The pairings in the two right-hand columns are induced by the pairing of
complexes of sheaves
\[
\symup{\Homsh_{\Osh_X}(L_{*},\underline{\Omega}_X^{p})
\times\Homsh_{\Osh_X}(L'_{*},\underline{\Omega}_X^{p'})
\longrightarrow
\Homsh_{\Osh_X}(L_{*}\otimes L'_{*},\underline{\Omega}_X^{p+p'})}
\]
defined using the exterior product
$\symup{\underline{\Omega}_X^{p}\times\underline{\Omega}_X^{p'}\longrightarrow\underline{\Omega}_X^{p+p'}}$.
The pairing in the first column is the cup product (relative to the exterior
product). I claim that the last row of (4.5) identifies with the diagram of
homomorphisms analogous to (4.3), with $\symup{Y}$ replaced by $\symup{Y\cap Y'}$ and $\symup{p}$ by
$\symup{p+p'}$. For this it suffices to show that $\symup{L\otimes L'}$ is an (evidently locally
free) resolution of $\symup{\Osh_{Y\cap Y'}}$. Indeed,
\[
  \symup{H_0(L\otimes L')=\Osh_Y\otimes\Osh_{Y'}=\Osh_{Y\cap Y'}}
\]
and
\[
  \symup{H_i(L\otimes L')=\operatorname{Tor}_{i}^{\Osh_X}(\Osh_Y,\Osh_{Y'})=0}
\]
for $\symup{i>0}$, because $\symup{Y}$ and $\symup{Y'}$ meet transversely. Theorem 1 now follows from the formula:
\begin{equation}\tag{4.6}
  \symup{s_Y.s_{Y'}=s_{Y.Y'}}
\end{equation}
(where the product on the left is that of the last column of (4.5)). Formula
(4.6), being purely local, is verified without difficulty by taking for $\symup{L_{*}}$
and $\symup{L'_{*}}$ the resolutions considered in Proposition 4. One proves similarly
(and more easily) that $\symup{Z\longrightarrow P_X(Z)}$ is compatible with cartesian
product:
\begin{equation}\tag{4.7}
  \symup{P_{X\times X'}(Z\times Z')=P_X(Z)\otimes P_{X'}(Z')}
\end{equation}
(this formula is valid when $\symup{Z}$, resp. $\symup{Z'}$, is a nonsingular cycle on
the nonsingular variety $\symup{X}$, resp. $\symup{X'}$, with $\symup{Z\times Z'}$ regarded as a
nonsingular cycle on $\symup{X\times X'}$). From (4.4) and (4.7), it follows that
$\symup{P_X(Z)}$ is also compatible with the operation \fquote{inverse image} under a morphism of nonsingular
varieties $\symup{f:X\longrightarrow X'}$:
\begin{equation}\tag{4.8}
  \symup{P_X\bigl(f^{-1}(Z')\bigr)=f^{*}\bigl(P_{X'}(Z')\bigr)}
\end{equation}
a formula valid if $\symup{Z'}$ is a nonsingular cycle on $\symup{X'}$ such that $\symup{f}$ is \fquote{transverse} to $\symup{Z'}$, i.e. such that the graph of $\symup{f}$ is transverse to the cycle $\symup{X\times Z'}$ in $\symup{X\times X'}$.

\src{181}
\result{COROLLARY}{1} \uline{Let} $\symup{X}$ \uline{and} $\symup{X'}$ \uline{be two nonsingular locally closed
varieties in a projective space; let us suppose} $\symup{X'}$ \uline{complete. Let} $\symup{U}$ \uline{be a nonsingular
cycle on} $\symup{X\times X'}$, \uline{and let} $\symup{a}$ \uline{and} $\symup{b}$ \uline{be two points of} $\symup{X'}$ \uline{such that} $\symup{U}$
\uline{meets the cycles} $\symup{X\times(a)}$ \uline{and} $\symup{X\times(b)}$ \uline{transversely; let} $\symup{Z}$ \uline{and} $\symup{Z'}$ \uline{be
the nonsingular cycles on} $\symup{X}$ \uline{such that}
$\symup{Z\times(a)=(X\times(a)).U}$, $\symup{Z'\times(b)=(X\times(b)).U}$.
\uline{Then one has}
\[
  \symup{P_X(Z)=P_X(Z').}
\]

Indeed, let $\symup{f_a:X\longrightarrow X\times X'}$ be defined by
$\symup{f_a(x)=(x,a)}$; one then has, by (4.8), the formula $\symup{P(Z)=f_a^{*}(P(U))}$, and similarly
$\symup{P(Z')=f_b^{*}(P(U))}$. Using the Künneth formula
\[
  \symup{H^{*}(X\times X',\underline{\Omega}_{X\times X'}^{*})
  =H^{*}(X,\underline{\Omega}_X^{*})\otimes H^{*}(X',\underline{\Omega}_{X'}^{*})}
\]
and the fact that $\symup{H^{0}(X',\Osh_{X'})}$ consists only of scalars, one readily
obtains $\symup{f_a^{*}=f_b^{*}}$, proving the assertion.

For every $\symup{x\in X}$, $\symup{(x)}$ is a nonsingular subvariety of $\symup{X}$ of
codimension $\symup{n}$, and therefore defines an element $\symup{\varepsilon_x}$ of
$\symup{H^{n}(X,\underline{\Omega}_X^{n})}$. If $\symup{X}$ is a nonsingular projective variety, Corollary 1
shows that $\symup{\varepsilon_x}$ does not depend on the chosen point $\symup{x}$; it is denoted
$\symup{\varepsilon_X}$ and called the \uline{fundamental class} of
$\symup{H^{n}(X,\underline{\Omega}_X^{n})}$.

\result{REMARK}{} To obtain a satisfactory theory, one would have to define
$\symup{P_X(Z)}$ for arbitrary cycles $\symup{Z}$ and prove Theorem 1 for a proper intersection
of cycles. (At the time of writing this exposé, this has not yet been done in
full generality). Assuming this has been done, Corollary 1 becomes: if $\symup{Z}$ and
$\symup{Z'}$ are two algebraically equivalent cycles, then $\symup{P_X(Z)=P_X(Z')}$ (a statement
that does not seem to follow from the preceding results, even when $\symup{Z}$ and $\symup{Z'}$
are nonsingular).

\fsection{\normalfont5}{\normalfont\uline{The duality theorem}.}

In this section, $\symup{X}$ denotes a nonsingular projective variety of dimension $\symup{n}$.

\result{THEOREM}{2} \uline{The fundamental class} $\symup{\varepsilon_X}$ \uline{of}
$\symup{H^{n}(X,\underline{\Omega}_X^{n})}$ \uline{is a basis of this vector space}.

(The proof will be given below). By the preceding theorem, one may therefore
identify $\symup{H^{n}(X,\underline{\Omega}_X^{n})}$ with the field $\symup{\underline{k}}$. Let us now consider the pairings
introduced in Section 2, which in particular give a pairing

\src{182}
\[
\symup{\operatorname{Ext}_{\Osh_X}^{p}(X;\Osh_X,F)
\times\operatorname{Ext}_{\Osh_X}^{n-p}(X;F,\underline{\Omega}_X^{n})
\longrightarrow
\operatorname{Ext}_{\Osh_X}^{n}(X;\Osh_X,\underline{\Omega}_X^{n}),\quad\text{i.e.}}
\]
\begin{equation}\tag{5.1}
  \symup{H^{p}(X,F)
  \times\operatorname{Ext}_{\Osh_X}^{n-p}(X;F,\underline{\Omega}_X^{n})
  \longrightarrow H^{n}(X,\underline{\Omega}_X^{n})}
\end{equation}
In view of Theorem 2, this pairing defines a homomorphism
\begin{equation}\tag{5.2}
  \symup{\operatorname{Ext}_{\Osh_X}^{n-p}(X;F,\underline{\Omega}_X^{n})
  \longrightarrow\bigl(H^{p}(X,F)\bigr)'}
\end{equation}
This homomorphism is functorial in $\symup{F}$ and commutes with the coboundary
homomorphisms associated with exact sequences
$\symup{0\longrightarrow F'\longrightarrow F
\longrightarrow F''\longrightarrow0}$.

\result{THEOREM}{3} \uline{The homomorphism} (5.2) \uline{is an isomorphism}.

In particular, one recovers Serre's result:

\result{COROLLARY}{} \uline{Let} $\symup{E}$ \uline{be an algebraic vector bundle on} $\symup{X}$, \uline{and let}
$\symup{\Osh_X(E)}$ \uline{be the sheaf of germs of regular sections of} $\symup{E}$; \uline{one then has
canonical isomorphisms}:
\begin{equation}\tag{5.3}
  \symup{\bigl(H^{p}(X,\Osh_X(E))\bigr)'
  =H^{n-p}\bigl(X,\underline{\Omega}_X^{n}\otimes\Osh_X(E')\bigr)}
\end{equation}
It suffices to apply Theorem 3 and Corollary 1 of Proposition 1.

Theorems 2 and 3 will follow from the following assertion:

\noindent(D) \uline{The following homomorphism}
\begin{equation}\tag{5.2 bis}
  \symup{\operatorname{Ext}_{\Osh_X}^{n-p}(X;F,\underline{\Omega}_X^{n})
  \longrightarrow\bigl(H^{p}(X,F)\bigr)'\otimes L
  \qquad (L=H^{n}(X,\underline{\Omega}_X^{n})),}
\end{equation}
\uline{induced by the pairing (5.1) is an isomorphism}.

Let us indeed show that (D) implies Theorem 2. Let $\symup{\underline{k}_x=\Osh_{(x)}}$ be the sheaf
of local rings of the variety reduced to the point $\symup{x\in X}$; let us consider the
canonical homomorphism $\symup{\Osh_X\longrightarrow \underline{k}_x}$ and the associated
homomorphism
\begin{equation}\tag{5.3 bis}
  \symup{H^{0}(X,\Osh_X)\longrightarrow H^{0}(X,\underline{k}_x)}
\end{equation}
its transpose identifies with the homomorphism
\begin{equation}\tag{5.4}
  \symup{\operatorname{Ext}_{\Osh_X}^{n}(X;\underline{k}_x,\underline{\Omega}_X^{n})\otimes L'
  \longrightarrow
  \operatorname{Ext}_{\Osh_X}^{n}(X;\Osh_X,\underline{\Omega}_X^{n})\otimes L'}
\end{equation}
induced by the homomorphism on $\symup{\operatorname{Ext}^{n}}$ associated with
$\symup{\Osh_X\longrightarrow \underline{k}_x}$:
\begin{equation}\tag{5.5}
  \symup{\operatorname{Ext}_{\Osh_X}^{n}(X;\underline{k}_x,\underline{\Omega}_X^{n})
  \longrightarrow
  \operatorname{Ext}_{\Osh_X}^{n}(X;\Osh_X,\underline{\Omega}_X^{n})
  =H^{n}(X,\underline{\Omega}_X^{n})}
\end{equation}

\src{183}
Since (5.3 bis) is an isomorphism, so is (5.4),
and hence so is (5.5). Since $\symup{s_{(x)}}$ is a basis of
$\symup{\operatorname{Ext}_{\Osh_X}^{n}(X;\underline{k}_x,\underline{\Omega}_X^{n})}$ by (4.2), it follows that
its image $\symup{\varepsilon_X}$ is a basis of $\symup{H^{n}(X,\underline{\Omega}_X^{n})}$.

It remains to prove assertion (D), which will follow in a purely formal way
from the elementary facts summarized in the following lemmas. Here $\symup{X}$ is
assumed to be a closed subset (singular or not) of projective space $\symup{P}$ of
dimension $\symup{r}$. One uses $\symup{\Osh_P(m)}$ for the sheaf on $\symup{P}$ denoted $\symup{\Osh(m)}$ in
[3], and $\symup{\Osh_X(m)}$ for the analogous sheaf on $\symup{X}$.

\result{LEMMA}{2} \uline{Assertion} (D) \uline{is true if} $\symup{X=P}$ \uline{and} $\symup{F=\Osh_P(m)}$

This lemma is verified by a direct calculation. The explicit calculation of
$\symup{H^{i}(P,\Osh_P(m))}$ is given in [3], but it can be carried out more
elementarily. The computation of the cup product
$\symup{H^{i}(P,\Osh_P(m))\times H^{j}(P,\Osh_P(m'))\longrightarrow H^{i+j}(P,\Osh_P(m+m'))}$
needed to make the pairing (5.1) explicit presents no
difficulty.

\result{LEMMA}{3} \uline{Every coherent algebraic sheaf} $\symup{F}$ \uline{on} $\symup{X}$ \uline{is isomorphic to a
quotient sheaf of a sheaf} $\symup{\Osh_X(-m)^{k}}$, \uline{where} $\symup{m}$ \uline{may be taken as large as
desired}.

This follows from the fact that $\symup{F\otimes\Osh_X(m)}$ is ``generated by its
sections'' for $\symup{m}$ large; cf. [3].

\result{LEMMA}{4} \uline{Let} $\symup{i>0}$. \uline{Then}
$\symup{H^{r-i}(P,\Osh_P(-m))=0}$ \uline{for} $\symup{m}$ \uline{large; and for every coherent algebraic sheaf}
$\symup{B}$ \uline{on} $\symup{X}$, \uline{one has}
$\symup{\operatorname{Ext}_{\Osh_X}^{i}(X;\Osh_X(-m),B)=0}$
\uline{for} $\symup{m}$ \uline{large}.

The first assertion follows from the explicit calculations mentioned above;
for the second, one notes that one has an isomorphism
\[
  \symup{\operatorname{Ext}_{\Osh_X}^{i}(X;\Osh_X(-m),B)
  =H^{i}(X,B\otimes\Osh(m))}
\]
(Proposition 1, Corollary 1), whence the conclusion by a well-known result of
[3]. Combining the two preceding lemmas, one obtains the

\result{COROLLARY}{} \uline{Let} $\symup{i>0}$. \uline{Then the functor}
$\symup{F\longrightarrow H^{r-i}(P,F)}$ \uline{on the category of coherent algebraic sheaves on}
$\symup{P}$ \uline{is coeffaceable, and so is the functor}
$\symup{\operatorname{Ext}_{\Osh_X}^{i}(X;F,B)}$ \uline{on the category of coherent algebraic
sheaves on} $\symup{X}$.

\src{184}
\result{LEMMA}{5} \uline{Let} $\symup{A}$ \uline{and} $\symup{B}$ \uline{be two coherent algebraic sheaves on} $\symup{X}$; \uline{let us
set} $\symup{A(m)=A\otimes\Osh_X(m)}$. \uline{Then, for} $\symup{m}$ \uline{sufficiently large, the canonical
homomorphism}
\[
  \symup{\operatorname{Ext}_{\Osh_X}^{i}(X;A(-m),B)
  \longrightarrow H^{0}\bigl(X,\Extsh_{\Osh_X}^{i}(A(-m),B)\bigr)
  =H^{0}\bigl(X,\Extsh_{\Osh_X}^{i}(A,B)(m)\bigr)}
\]
\uline{is an isomorphism}.

This follows at once from the spectral sequence of Proposition 1 applied to
$\symup{A(-m)}$ and $\symup{B}$, since one then has
\[
  \symup{E_{2}^{p,q}(A(-m),B)
  =H^{p}\bigl(X,\Extsh_{\Osh_X}^{q}(A(-m),B)\bigr)
  =H^{p}\bigl(X,\Extsh_{\Osh_X}^{q}(A,B)(m)\bigr),}
\]
which vanishes for $\symup{p>0}$ and $\symup{m}$ large.

Let us now prove (D) in the case $\symup{X=P}$. First we prove that (5.2 bis) is an
isomorphism for $\symup{p=n}$; since both sides are then left-exact functors (because
$\symup{H^{r+1}(P,F)=0}$), Lemma 3 shows that it suffices to prove the assertion for
$\symup{F=\Osh_P(-m)}$; in that case it is contained in Lemma 2. Since the homomorphisms
(5.2 bis) are functorial and compatible with the coboundary operators, and
for $\symup{p<n}$ its two sides are coeffaceable functors of $\symup{F}$ (Corollary to Lemma 4),
a standard argument therefore shows that (5.2 bis) is an isomorphism for every
$\symup{p}$. This proves the duality theorem for projective space.

Let us now suppose $\symup{X}$ arbitrary but nonsingular. By the duality theorem for $\symup{P}$, one has an
isomorphism,
\[
  \symup{H^{n}(X,F)'=H^{n}(P,F)'
  =\operatorname{Ext}_{\Osh_P}^{r-n}(P;F,\underline{\Omega}_P^{r})}
\]
By Lemma 1 (no. 4), the last term identifies with
\[
  \symup{\operatorname{Hom}_{\Osh_P}(P;F,\underline{\omega}_X^{n})
  =\operatorname{Hom}_{\Osh_X}(X;F,\underline{\Omega}_X^{n})
  =\operatorname{Ext}_{\Osh_X}^{0}(X;F,\underline{\Omega}_X^{n})}
\]
and hence there is an isomorphism
\begin{equation}\tag{5.6}
  \symup{H^{n}(X,F)'=\operatorname{Hom}_{\Osh_X}(X;F,\underline{\Omega}_X^{n})
  =\operatorname{Ext}_{\Osh_X}^{0}(X;F,\underline{\Omega}_X^{n})}
\end{equation}
Taking $\symup{F=\underline{\Omega}_X^{n}}$, one obtains an isomorphism
\begin{equation}\tag{5.7}
  \symup{\gamma:H^{n}(X,\underline{\Omega}_X^{n})\xrightarrow{\sim}\underline{k}}
\end{equation}

\src{185}
One checks that the isomorphism (5.6) is precisely (5.2 bis) for $\symup{p=n}$ after
setting $\symup{L=\underline{k}}$ by means of (5.7). Consequently, (5.2 bis) is an isomorphism for
$\symup{p=n}$. To prove that it is an isomorphism for every $\symup{p}$, it again suffices to
prove that, for $\symup{p<n}$, the two sides of (5.2 bis) are coeffaceable functors of
$\symup{F}$, and a fortiori (in view of Lemma 3) that both sides vanish when
$\symup{F=\Osh_X(-m)}$ with $\symup{m}$ large. For the first side this follows from Lemma 4; for
the second, using the duality theorem on $\symup{P}$, one writes:
\[
  \symup{H^{p}(X,\Osh_X(-m))'
  =\operatorname{Ext}_{\Osh_P}^{r-p}
    (P;\Osh_X(-m),\underline{\Omega}_P^{r}).}
\]
The second side vanishes for $\symup{p<n}$ and $\symup{m}$ large, as follows from Lemma 5 (with
$\symup{X=P}$ there) and from the fact that, as a coherent algebraic sheaf on $\symup{P}$,
$\symup{\Osh_X}$ has cohomological dimension $\symup{\leq r-n}$ (since $\symup{X}$ is nonsingular),
whence
\[
  \symup{\Extsh_{\Osh_P}^{r-p}(\Osh_X,\underline{\Omega}_P^{r})=0
  \qquad\text{for }p<n.}
\]

\fsection{\normalfont6}{\normalfont\uline{The duality theorem for singular varieties}.}

Let $\symup{X}$ be a closed subset of dimension $\symup{n}$ of projective space $\symup{P}$ of
dimension $\symup{r}$. Formula (5.6) must here be written
\begin{equation}\tag{6.1}
  \symup{H^{n}(X,F)'
  \simeq\operatorname{Hom}_{\Osh_X}(X;F,\underline{\omega}_X^{n})
  =\operatorname{Ext}_{\Osh_X}^{0}(X;F,\underline{\omega}_X^{n})}
\end{equation}
where one has set
\begin{equation}\tag{6.2}
  \symup{\underline{\omega}_X^{n}=E^{0}(\Osh_X)
  =\Extsh_{\Osh_P}^{r-n}(\Osh_X,\underline{\Omega}_P^{r})}
\end{equation}
As stated in Proposition 5, the sheaf so introduced does not in fact depend on
the chosen immersion of $\symup{X}$ into a nonsingular variety $\symup{P}$. Taking
$\symup{F=\underline{\omega}_X^{n}}$ in (6.1), one obtains
\begin{equation}\tag{6.2 bis}
  \symup{H^{n}(X,\underline{\omega}_X^{n})'
  \simeq\operatorname{Hom}_{\Osh_X}(X;\underline{\omega}_X^{n},\underline{\omega}_X^{n})}
\end{equation}
and hence a distinguished element
of $\symup{H^{n}(X,\underline{\omega}_X^{n})'}$ corresponding to the identity homomorphism of
$\symup{\underline{\omega}_X^{n}}$ into itself:
\begin{equation}\tag{6.3}
  \symup{\gamma:H^{n}(X,\underline{\omega}_X^{n})\longrightarrow\underline{k}}
\end{equation}

\src{186}
Let us now consider the pairings defined by composition of $\symup{\operatorname{Ext}}$:
\begin{equation}\tag{6.4}
  \symup{H^{p}(X,F)\times
  \operatorname{Ext}_{\Osh_X}^{n-p}(X;F,\underline{\omega}_X^{n})
  \longrightarrow H^{n}(X,\underline{\omega}_X^{n})}
\end{equation}
and let us compose them with the homomorphism $\symup{\gamma}$ of (6.3); one obtains functorial
homomorphisms, compatible with the boundary operators:
\begin{equation}\tag{6.5}
  \symup{\operatorname{Ext}_{\Osh_X}^{n-p}(X;F,\underline{\omega}_X^{n})
  \longrightarrow H^{p}(X,F)'}
\end{equation}
(generalizing (5.2)). One checks that for $\symup{p=n}$ this gives the isomorphism
(6.1). With this understood, one has the

\result{THEOREM}{3 bis} \uline{The following four conditions on} $\symup{X}$ \uline{are equivalent for
a given integer} $\symup{k\geq0}$:
\begin{enumerate}
  \item[\textnormal{i.}] \uline{The functorial homomorphism} (6.5) \uline{is an isomorphism for}
    $\symup{n-k\leq p\leq n}$.
  \item[\textnormal{ii.}] \uline{One has} $\symup{H^{p}(X,\Osh_X(-m))=0}$ \uline{for} $\symup{m}$ \uline{large and}
    $\symup{n-k\leq p<n}$.
  \item[\textnormal{iii.}] \uline{The functor} $\symup{H^{p}(X,F)}$ \uline{on the category of coherent
    algebraic sheaves on} $\symup{X}$ \uline{is coeffaceable for} $\symup{n-k\leq p<n}$.
  \item[\textnormal{iv.}] \uline{One has}
    $\symup{E^{i}(\Osh_X)=\Extsh_{\Osh_P}^{r-n+i}(\Osh_X,\underline{\omega}_P^{r})=0}$
    \uline{for} $\symup{0<i\leq k}$.
\end{enumerate}

\noindent PROOF. -- \textnormal{i.}$\symup{\Rightarrow}$\textnormal{ii.} follows
from Lemma 4, \textnormal{ii.}$\symup{\Rightarrow}$\textnormal{iii.} from Lemma 3, and
\textnormal{iii.}$\symup{\Rightarrow}$\textnormal{i.} by a well-known standard argument, since
the two sides of (6.5) are then coeffaceable functors for $\symup{n-k\leq p<n}$ (the
first by Lemma 4). Finally, let us prove
\textnormal{ii.}$\symup{\Leftrightarrow}$\textnormal{iv.} This follows from the corollary to the
next proposition:

\result{PROPOSITION}{6} \uline{Let} $\symup{F}$ \uline{be a coherent algebraic sheaf on} $\symup{X}$, \uline{and let}
$\symup{i}$ \uline{be an integer. Then, for} $\symup{m}$ \uline{sufficiently large, one has an isomorphism}:
\begin{equation}\tag{6.6}
  \symup{H^{i}(X,F(-m))'
  \simeq H^{0}\bigl(X,E^{n-i}(F)(m)\bigr),}
\end{equation}
\uline{where one sets} (compare no. 3, Proposition 5):
\begin{equation}\tag{6.7}
  \symup{E^{j}(F)=\Extsh_{\Osh_P}^{r-n+j}(F,\underline{\Omega}_P^{r})}
\end{equation}
Indeed, by the duality theorem on $\symup{P}$, the first side of (6.6) is isomorphic to
$\symup{\operatorname{Ext}_{\Osh_P}^{r-i}(P;F(-m),\underline{\Omega}_P^{r})}$, so (6.6) follows
from Lemma 5.

\src{187}
\result{COROLLARY}{} \uline{For} $\symup{H^{i}(X,F(-m))=0}$ \uline{to hold for} $\symup{m}$ \uline{large, it is necessary and sufficient that}
$\symup{E^{n-i}(F)=0}$.

Let us recall that the $\symup{E^{j}(F)}$ do not in fact depend on the projective immersion
under consideration. The condition in the corollary is purely local;
consequently, if it holds for $\symup{F}$, it holds for every sheaf locally isomorphic
to a sheaf $\symup{F^{n}}$. In particular, if this condition holds for $\symup{\Osh_X}$, it
holds for every locally free coherent algebraic sheaf. This is the case, for
example, for every $\symup{i<n}$ if $\symup{X}$ is nonsingular; for $\symup{i=0}$ if no component of
$\symup{X}$ is reduced to a single point; and for $\symup{i=0,1}$ if $\symup{X}$ is normal and all its
components have dimension $\symup{>1}$ (cf. [3]). For this to
hold for every $\symup{i<k}$, it is necessary and sufficient, by definition, that the
local rings $\symup{\Osh_x}$ ($\symup{x\in X}$) have ``homological codimension'' $\symup{\geq k}$ (cf.
[4] for this notion). If $\symup{k=n}$, Theorem 3 bis says that this means the duality
theorem is true for $\symup{X}$, i.e. that (6.5) is an isomorphism for every $\symup{p}$ and
every $\symup{F}$. Many equivalent conditions on the local rings $\symup{\Osh_x}$ can be given
for this to hold (NAGATA), for example that they satisfy the Cohen-Macaulay
equidimensionality theorem. It holds, for example, if $\symup{X}$ is locally a
``complete intersection'' in a nonsingular ambient variety.

\fsection{\normalfont7}{\normalfont\uline{Poincaré duality}.}

Let $\symup{X}$ be a nonsingular projective variety of dimension $\symup{n}$. Then
$\symup{H^{**}(X)=H^{*}(X,\underline{\Omega}_X^{*})}$
is a finite-dimensional anticommutative bigraded algebra, which we grade by
total degree, so that $\symup{H^{p,q}(X)=H^{p}(X,\underline{\Omega}_X^{q})}$ has degree $\symup{p+q}$. The
degrees of $\symup{H^{*}(X)}$ range from $\symup{0}$ to $\symup{2n}$. By Theorem 2 and the corollary to
Theorem 3, $\symup{H^{*}(X)}$ is a ``Poincaré algebra'' of dimension $\symup{2n}$, i.e.
$\symup{H^{2n}(X)}$ is equipped with an isomorphism to the base field $\symup{\underline{k}}$, and the
product
$\symup{H^{m}(X)\times H^{2n-m}(X)\longrightarrow H^{2n}(X)=\underline{k}}$
is a duality between $\symup{H^{m}(X)}$ and $\symup{H^{2n-m}(X)}$. Moreover, if $\symup{Y}$ is another
nonsingular projective variety, the Künneth formula for coherent algebraic
sheaves gives
\begin{equation}\tag{7.1}
  \symup{H^{*}(X\times Y)=H^{*}(X)\otimes H^{*}(Y)}
\end{equation}
an isomorphism compatible with the Poincaré-algebra structures. Furthermore,
$\symup{H^{*}(X)}$ is, as a commutative algebra, a contravariant functor of $\symup{X}$, a
morphism $\symup{f:Y\longrightarrow X}$ defining in the evident way a homomorphism of
graded algebras

\src{188}
\begin{equation}\tag{7.2}
  \symup{f^{*}:H^{*}(X)\longrightarrow H^{*}(Y)}
\end{equation}
Since these are Poincaré algebras, one obtains by transposition a homomorphism of vector
spaces:
\begin{equation}\tag{7.3}
  \symup{f_{*}:H^{*}(Y)\longrightarrow H^{*}(X)}
\end{equation}
One has seen in no. 4 that the effect of $\symup{f^{*}}$ on the cohomology classes
corresponding to nonsingular cycles has the geometric interpretation of taking
the cohomology classes corresponding to their inverse images. In the present
setting it is important to show that (7.3) similarly corresponds to the
``direct image'' operation on cycles. Under suitable nonsingularity conditions,
at least, this follows from the following special case:

\result{THEOREM}{4} \uline{If} $\symup{f}$ \uline{is the inclusion of a nonsingular subvariety} $\symup{Y^{m}}$
\uline{of} $\symup{X^{n}}$ \uline{into} $\symup{X^{n}}$, \uline{then, writing} $\symup{1_Y}$ \uline{for the unit element of} $\symup{H(Y)}$,
\uline{one has}
\begin{equation}\tag{7.4}
  \symup{f_{*}(1_Y)=P_X(Y)}
\end{equation}
\uline{where the right-hand side is the cohomology class in} $\symup{X}$ \uline{associated with} $\symup{Y}$.

This formula is equivalent to
\begin{equation}\tag{7.4 bis}
  \symup{\bigl\langle\xi^{m,m},P_X(Y)\bigr\rangle\varepsilon_Y
  =f^{*}(\xi^{m,m})
  \qquad
  \bigl(\xi^{m,m}\in H^{m}(X,\underline{\Omega}_X^{m})\bigr)}
\end{equation}
 where $\symup{\varepsilon_Y}$ is
the fundamental element in $\symup{H^{m}(Y,\underline{\Omega}_Y^{m})}$, and, for nonsingular projective
varieties, this formula gives a new definition of the cohomology class associated with
$\symup{Y}$. To prove Theorem 4, one interprets, using Theorem 3, the transpose of the
homomorphism
\[
  \symup{H^{m}(X,\underline{\Omega}_X^{m})
  \longrightarrow H^{m}(Y,\underline{\Omega}_Y^{m})
  =H^{m}(X,\underline{\Omega}_Y^{m})}
\]
as the homomorphism
\begin{equation}\tag{7.5}
\symup{H^{n-m}(X,\underline{\Omega}_X^{n-m})
  \xrightarrow{\sim}
  \operatorname{Ext}_{\Osh_X}^{n-m}(X;\underline{\Omega}_X^{m},\underline{\Omega}_X^{n})
  \longleftarrow
  \operatorname{Ext}_{\Osh_X}^{n-m}(X;\underline{\Omega}_Y^{m},\underline{\Omega}_X^{n})
  \xrightarrow{\sim}
  \operatorname{Hom}_{\Osh_X}(X;\underline{\Omega}_Y^{m},\underline{\Omega}_Y^{m})}
\end{equation}
One checks that the element $\symup{1_Y}$ of the dual of
$\symup{H^{m}(Y,\underline{\Omega}_Y^{m})}$ identifies with the element of the right-hand side
corresponding to the identity endomorphism of $\symup{\underline{\Omega}_Y^{m}}$, and also that the
image of this element in $\symup{H^{n-m}(X,\underline{\Omega}_X^{n-m})}$ is indeed $\symup{P_X(Y)}$. These
compatibilities, which could have been stated already in no. 4, can be
formulated, and in fact hold, for arbitrary nonsingular varieties, the second

\src{189}
following, for example, from the commutativity of the following diagram of
canonical homomorphisms:
\begin{equation}\tag{7.6}
\resizebox{\textwidth}{!}{$
\left\{
\begin{array}{ccccc}
\symup{\operatorname{Ext}_{\Osh_X}^{n-m}(X;\underline{\Omega}_X^{m},\underline{\Omega}_X^{n})}
&
\symup{\longleftarrow}
&
\symup{\operatorname{Ext}_{\Osh_X}^{n-m}(X;\underline{\Omega}_Y^{m},\underline{\Omega}_X^{n})}
&
\symup{\xrightarrow{\sim}}
&
\symup{\operatorname{Hom}_{\Osh_X}(X;\underline{\Omega}_Y^{m},\underline{\Omega}_Y^{m})}
\\[3pt]
\symup{\Big\downarrow}
&&&&
\symup{\Big\downarrow}
\\[3pt]
\symup{\operatorname{Ext}_{\Osh_X}^{n-m}(X;\Osh_X,\underline{\Omega}_X^{n-m})}
&
\symup{\longleftarrow}
&
\symup{\operatorname{Ext}_{\Osh_X}^{n-m}(X;\Osh_Y,\underline{\Omega}_X^{n-m})}
&
\symup{\xrightarrow{\sim}}
&
\symup{\operatorname{Hom}_{\Osh_X}(X;\underline{\Omega}_X^{m},\underline{\Omega}_Y^{m})}
\end{array}
\right.
$}
\end{equation}

One thus obtains an exact analogue of the formalism of Poincaré duality on compact
oriented manifolds. In particular, Theorem 4 determines the cohomology class
associated with the diagonal of $\symup{X\times X}$. By a well-known argument, one deduces from this,
for example, a \uline{Lefschetz formula}:

\result{THEOREM}{5} \uline{Let} $\symup{f}$ \uline{be an endomorphism of a nonsingular projective
variety} $\symup{X}$, \uline{such that the fixed points of} $\symup{f}$ \uline{have multiplicity} $\symup{1}$. \uline{Then the number of
these fixed points is congruent}, $\symup{\operatorname{mod}}$ \uline{the characteristic of} $\symup{\underline{k}}$ \uline{to the
alternating sum of the traces of the endomorphisms of the} $\symup{H^{i}(X)}$ \uline{induced by}
$\symup{f}$.

The restriction imposed on $\symup{f}$ is due to the difficulties mentioned in the
remark in no. 4. One will note, however, that the Lefschetz formula remains true if
$\symup{f}$ is ``homotopic'' to an endomorphism all of whose fixed points have
multiplicity $\symup{1}$.

\fsection{\normalfont8}{\normalfont\uline{Generalization of the duality theorem}.}

Let $\symup{X}$ be a nonsingular algebraic variety such that every coherent algebraic
sheaf $\symup{F}$ on $\symup{X}$ is isomorphic to a quotient of a locally free coherent
algebraic sheaf (as is the case if $\symup{X}$ is locally closed in projective space).
Then every coherent algebraic sheaf $\symup{F}$ on $\symup{X}$ admits a finite resolution $\symup{L}$
by locally free sheaves; for any two such resolutions one can always find a
third and a homomorphism from it to each of the first two compatible with the
augmentations. Likewise, if $\symup{L}$ is a finite locally free resolution of $\symup{F}$ and
a homomorphism $\symup{F'\longrightarrow F}$ is given, there exists a finite locally
free resolution $\symup{L'}$ of $\symup{F'}$ and a homomorphism $\symup{L'\longrightarrow L}$
compatible with $\symup{F'\longrightarrow F}$, which may be assumed surjective if
$\symup{F'\longrightarrow F}$ is surjective. Given two integers $\symup{r,s\geq0}$, this makes
it possible to define two cohomological multifunctors in the arguments
$\symup{A_1,\ldots,A_r}$; $\symup{B_1,\ldots,B_s}$

\src{190}
which range over the category of coherent algebraic sheaves on $\symup{X}$ and take
values in the category of coherent algebraic sheaves on $\symup{X}$, resp.
in the category of modules over $\symup{H^{0}(X,\Osh_X)}$, by the formulas:
\begin{equation}\tag{8.1}
\resizebox{\textwidth}{!}{$
\left\{
\begin{aligned}
\symup{\underline{T}_{r}^{s*}(A_1,\ldots,A_r;B_1,\ldots,B_s)}
&
\symup{={}H^{*}\!\left(
    \Homsh_{\Osh_X}\!\left(
      \underline{L}(A_1)\otimes\cdots\otimes \underline{L}(A_r),
      \underline{L}(B_1)\otimes\cdots\otimes \underline{L}(B_s)
    \right)\right)}
\\[3pt]
\symup{T_{r}^{s*}(A_1,\ldots,A_r;B_1,\ldots,B_s)}
&
\symup{={}\underline{\underline{R}}^{*}\Gamma\!\left(
    \Homsh_{\Osh_X}\!\left(
      \underline{L}(A_1)\otimes\cdots\otimes \underline{L}(A_r),
      \underline{L}(B_1)\otimes\cdots\otimes \underline{L}(B_s)
    \right)\right)}
\end{aligned}
\right.
$}
\end{equation}
Here $\symup{\underline{L}(F)}$ denotes a finite locally free resolution of the coherent algebraic
sheaf $\symup{F}$, and $\symup{\underline{\underline{R}}^{*}\Gamma(K)}$ denotes the hypercohomology of $\symup{X}$ with
respect to the complex of sheaves $\symup{K}$. If $\symup{r}$ or $\symup{s}$ is zero, the tensor product
of the $\symup{\underline{L}(A_i)}$, resp. of the $\symup{\underline{L}(B_j)}$, is replaced by $\symup{\Osh_X}$. In
particular, $\symup{\underline{T}_{0}^{0}}$ and $\symup{T_{0}^{0}}$ are graded functors with no
arguments, $\symup{\underline{T}_{0}^{0}}$ is concentrated in degree 0, where it is the
sheaf $\symup{\Osh_X}$, while $\symup{T_{0}^{0}}$ is identified with $\symup{H^{*}(X,\Osh_X)}$. The
independence of the right-hand sides of (8.1) from the chosen resolutions is in
any event evident for the first line (the question then being local), and for
the second follows from the preceding general remarks, taking into account the
hypercohomology spectral sequence of the complex of sheaves
$\symup{K=\Homsh_{\Osh_X}\bigl(\underline{L}(A_1)\otimes\cdots,
      \underline{L}(B_1)\otimes\cdots\bigr)}$,
which abuts to the hypercohomology of $\symup{X}$ with respect to $\symup{K}$, and whose initial term
is $\symup{H^{p}(X,H^{q}(K))}$, i.e.
\begin{equation}\tag{8.2}
  \symup{E_{2}^{p,q}
  =H^{p}\!\left(X,
    (\underline{T}_{r}^{s})^{(q)}
      (A_1,\ldots,A_r;B_1,\ldots,B_s)\right)}
\end{equation}
One then sees that this spectral sequence itself is independent of the chosen
resolutions, and its abutment is the first member of (8.1). One readily defines the coboundary
operators in the various arguments $\symup{A_i}$, $\symup{B_j}$, noting that every
exact sequence
$\symup{0\longrightarrow F'\longrightarrow F\longrightarrow F''\longrightarrow0}$
can be resolved by an exact sequence of finite locally free complexes.

On the system of functors $\symup{\underline{T}_{r}^{s*}}$, resp. $\symup{T_{r}^{s*}}$,
one defines operations analogous to those of tensor calculus, whose definitions
follow immediately from (8.1). Thus one has a \uline{composition} (generalizing the
one considered in no. 2):
\begin{equation}\tag{8.3}
\resizebox{\textwidth}{!}{$\symup{
T_{r}^{s*}(A_1,\ldots,A_r;B_1,\ldots,B_s)
\times T_{r'}^{s'*}(A'_1,\ldots,A'_{r'};B'_1,\ldots,B'_{s'})
\longrightarrow T_{r+r'}^{s+s'*}
(A_1,\ldots,A'_{r'};B_1,\ldots,B'_{s'})
}$}
\end{equation}
satisfying the evident associativity and compatibility properties with respect
to functorial homomorphisms, coboundary homomorphisms, and spectral sequences.
Likewise, one has symmetry operations, whose explicit description is left to
the reader. In addition, one has a \uline{contraction} operation whenever one of the
arguments $\symup{A_i}$

\src{191}
is equal to one of the arguments $\symup{B_j}$;
\begin{equation}\tag{8.4}
\resizebox{\textwidth}{!}{$
\begin{aligned}
&
\symup{T_{r}^{s*}(A_1,\ldots,A_{i-1},C,A_{i+1},\ldots,A_r;
 B_1,\ldots,B_{j-1},C,B_{j+1},\ldots,B_s)\longrightarrow}
\\
&
\symup{\qquad\relbar\quad
T_{r-1}^{s-1*}(A_1,\ldots,\widehat{A}_i,\ldots,A_r;
 B_1,\ldots,\widehat{B}_j\ldots,B_s)}
\end{aligned}
$}
\end{equation}
Furthermore, if an argument $\symup{A_i}$ is a locally free sheaf, then one may remove it
provided that one of the $\symup{A_j}$ ($\symup{j\neq i}$) is replaced by $\symup{A_j\otimes A_i}$, or
one of the $\symup{B_k}$ by $\symup{B_k\otimes A_i'}$ (where
$\symup{A_i'=\Homsh_{\Osh_X}(A_i,\Osh_X)}$); one has an analogous calculation rule
when one of the arguments $\symup{B_j}$ is locally free. In particular, one may always remove an argument
equal to $\symup{\Osh_X}$. If all the arguments are locally free
except for at most one of the $\symup{B_i}$, the rule just stated gives a functorial
isomorphism
\begin{equation}\tag{8.5}
  \symup{T_{r}^{s*}(A_1,\ldots,A_r;B_1,\ldots,B_s)
  =H^{*}\bigl(X,A_1'\otimes\cdots\otimes A_r'
    \otimes B_1\otimes\cdots\otimes B_s\bigr)}
\end{equation}
(since this reduces to the case $\symup{r=0}$, $\symup{s=1}$, where it is immediate: one may
also use directly the spectral sequence with initial term (8.2)). Operations
corresponding to the preceding ones are defined for the
$\symup{\underline{T}_{r}^{s}}$. The relations among the various operations so
introduced are the same as for the analogous operations in tensor calculus.

Let $\symup{n}$ be the dimension of $\symup{X}$. By successively applying a tensor composition
(8.3) and contractions (8.4) to the repeated arguments, one obtains a pairing
\begin{equation}\tag{8.6}
\resizebox{\textwidth}{!}{$\symup{
  (T_{r}^{s})^{p}(A_1,\ldots;B_1,\ldots)
  \times
  (T_{s}^{r})^{n-p}(B_1,\ldots;A_1,\ldots,A_r\otimes\underline{\Omega}_X^{n})
  \longrightarrow H^{n}(X,\underline{\Omega}_X^{n})
}$}
\end{equation}

\result{THEOREM}{6} \uline{If} $\symup{X}$ \uline{is a nonsingular projective variety, then the
pairings} (8.6) \uline{are dualities}.

This follows in a purely formal way from the corollary to Theorem 3. Indeed, it
follows readily from that corollary that if $\symup{K}$ is a complex of \uline{locally free}
coherent algebraic sheaves, then the hypercohomology of $\symup{X}$ with respect to
$\symup{K}$ is in duality with the hypercohomology of $\symup{X}$ with respect to
$\symup{K'\otimes\underline{\Omega}_X^{n}}$ through the natural pairings
\begin{equation}\tag{8.7}
  \symup{\underline{\underline{R}}^{p}\Gamma(K)
  \times\underline{\underline{R}}^{n-p}\Gamma(K'\otimes\underline{\Omega}_X^{n})
  \longrightarrow\underline{\underline{R}}^{n}\Gamma(\underline{\Omega}_X^{n})
  =H^{n}(X,\underline{\Omega}_X^{n})}
\end{equation}
One sees this by using the spectral sequence with initial term
$\symup{H^{p}(H^{q}(X,K))}$ and the analogous spectral sequence for
$\symup{K'\otimes\underline{\Omega}_X^{n}}$. Theorem 6 follows from the preceding result by using
definition (8.1).

\src{192}
\subsection*{\normalfont REMARKS.}

\noindent 1\textsuperscript{o} For the definitions preceding Theorem 6, it was not necessary for
$\symup{X}$ to be nonsingular, since it was not essential to work only with \uline{finite}
resolutions. But if $\symup{X}$ is singular, one can no longer assert a priori that the
$\symup{(\underline{T}_{r}^{s})^{p}(A_1,\ldots;B_1,\ldots)}$ are \uline{coherent} sheaves,
because the complex of sheaves
\[
  \symup{\Homsh_{\Osh_X}\bigl(\underline{L}(A_1)\otimes\cdots,\underline{L}(B_1)\otimes\cdots\bigr)}
\]
will have infinitely many components of a given total degree.

\medskip
\noindent 2\textsuperscript{o} One readily checks that in the formulas (8.1), one may replace \uline{one} of the
$\symup{\underline{L}(B_i)}$ by $\symup{B_i}$. In view of Proposition 3, this therefore shows that one has
\begin{equation}\tag{8.8}
\left\{
\begin{aligned}\symup{\underline{T}_{1}^{1*}(A;B)}&\symup{=\Extsh_{\Osh_X}^{*}(A,B)}\\\symup{T_{1}^{1*}(A;B)}&\symup{=\operatorname{Ext}_{\Osh_X}^{*}(X;A,B)}\end{aligned}
\right.
\end{equation}
In particular, taking $\symup{r=s=1}$ and $\symup{A_1=\Osh_X}$ in (8.6) recovers Theorem 3.
Formula (8.8) also implies
$\symup{T_{0}^{1*}(B)=H^{*}(X,B)}$, $\symup{T_{1}^{0*}(A)=\operatorname{Ext}_{\Osh_X}^{*}(X;A,\Osh_X)}$

\medskip
\noindent 3\textsuperscript{o} Formula (8.1) shows that in general the functors
$\symup{\underline{T}_{r}^{s*}}$ and $\symup{T_{r}^{s*}}$ have components in both positive \uline{and}
negative degrees. Using Remark 2, one sees that if $\symup{X}$ has dimension $\symup{n}$, then
the nonzero components of $\symup{\underline{T}_{r}^{s*}}$ lie between $\symup{-(s-1)n}$ and
$\symup{rn}$ if $\symup{s>0}$, and between $\symup{0}$ and $\symup{rn}$ if $\symup{s=0}$; the nonzero components of
$\symup{T_{r}^{s*}}$ lie between $\symup{-(s-1)n}$ and $\symup{(r+1)n}$ if $\symup{s>0}$, and between $\symup{0}$ and
$\symup{(r+1)n}$ if $\symup{s=0}$ (and, barring error, if $\symup{r>0}$, even between
$\symup{-(s-1)n}$ and $\symup{rn}$, resp. between $\symup{0}$ and $\symup{rn}$).

\section*{\centering\normalfont BIBLIOGRAPHY}

\noindent[1] CARTIER (Pierre). - Les groupes
$\symup{\operatorname{Ext}^{s}(A,B)}$, Séminaire A. Grothendieck : Algèbre homologique,
t. 1, 1957, n\textsuperscript{o} 3.

\noindent[2] GROTHENDIECK (Alexandre). - Sur quelques points d'algèbre
homologique, Tohoku math. J., t. 9, 1957, p. 119-183.

\noindent[3] SERRE (Jean-Pierre). - Faisceaux algébriques cohérents,
Annals of Math., t. 61, 1955, p. 197-278.

\noindent[4] SERRE (Jean-Pierre). - Sur la dimension homologique des anneaux
et des modules noethériens, Proc. Intern. Symp. on alg. number Theory
[1955. Tokyo et Nikko] - Tokyo, Science Council of Japan, 1956 ;
p. 175-189.

\clearpage
\src{193}
\section*{\centering\normalfont ADDENDUM}

The difficulties noted in the remark on page 13 have now been completely
resolved, through an extension of the duality theorems to arbitrary varieties
(with arbitrary singularities). For this theory, which can be formulated in
the general framework of the “theory of schemes”, we refer to the
“Eléments de Géométrie algébrique” of J. DIEUDONNÉ and A. GROTHENDIECK,
currently in preparation. One will find some indications in:

\begin{quote}
GROTHENDIECK (A.). - Cohomology theory of algebraic varieties,
Int. Congress of Mathematicians, 1958, Edinburgh (to appear).
\end{quote}

\begin{flushright}
[April 1959]
\end{flushright}

\begin{center}\rule{36mm}{0.5pt}\end{center}

\par\smallskip
\begingroup\footnotesize
\endgroup
